\input{preamble}

\title{Section 1.1 --- Sets and Numbers --- Answer Key}

\newlength{\lwid}
\newlength{\lht}

\newcommand{\Lcirc}[1]{%
\setlength{\lwid}{\widthof{$\mathbb{#1}$}}
\setlength{\lht}{\heightof{$\mathbb{#1}$}}
\begin{tikzpicture}[baseline=0,actualtext=#1]
\useasboundingbox (-0.5\lwid,0) rectangle (0.5\lwid,\lht);
\draw (0,0) node[anchor=base]{$\mathbb{#1}$};
\draw (0,0.5\lht) circle (1.75ex);
\end{tikzpicture}%
}

\newcommand{\sets}{\ensuremath{\mathbb{N}\quad\mathbb{Z}\quad\mathbb{Q}\quad\mathbb{R}}}
\newcommand{\setsN}{\ensuremath{\Lcirc{N}\quad\Lcirc{Z}\quad\Lcirc{Q}\quad\Lcirc{R}}}
\newcommand{\setsZ}{\N\tquad\Lcirc{Z}\tquad\Lcirc{Q}\tquad\Lcirc{R}}
\newcommand{\setsQ}{\ensuremath{\N\quad\Z\quad\Lcirc{Q}\quad\Lcirc{R}}}
\newcommand{\setsR}{\N\tquad\Z\tquad\Q\tquad\Lcirc{R}}

\newcommand{\picksets}[2][\sets]{$#2\hfill#1\qquad$}

\begin{document}
\maketitle

\section*{Standard Number Sets}
For each number below, circle the symbol for each set the number belongs to.
\begin{smenum}
\mitemxxx
    {\picksets[\setsR]{-\sqrt{7}}}
    {\picksets[\setsZ]{-\frac{75}{5}}}
    {\picksets[\setsQ]{\frac{1}{5}}}
\mitemxxx
    {\picksets[\setsZ]{-6}}
    {\picksets[\setsN]{5}}
    {\picksets[\setsN]{2}}
\mitemxxx
    {\picksets[\setsQ]{\frac{7}{3}}}
    {\picksets[\setsN]{\frac{21}{3}}}
    {\picksets[\setsQ]{3.61}}
\mitemxxx
    {\picksets[\setsR]{\pi}}
    {\picksets[\setsZ]{-4.0}}
    {\picksets[\setsZ]{-4}}
\mitemxxx
    {\picksets[\setsR]{\sqrt{8}}}
    {\picksets[\setsN]{\sqrt{9}}}
    {\picksets[\setsQ]{-4.\overline{54}}}
\mitemxxx
    {\picksets[\setsQ]{-0.2}}
    {\picksets[\setsZ]{-3215}}
    {\picksets[\setsN]{9.0}}
\mitemxxx
    {\picksets[\setsN]{816}}
    {\picksets[\setsQ]{-\frac{10}{7}}}
    {\picksets[\setsQ]{0.\overline{123}}}
\mitemxxo
    {\picksets[\setsQ]{-\frac{3}{4}}}
    {$0.1001110000...\qquad\setsR$}
\end{smenum}
\emph{Note: Answers here will vary. Those given are just some possible examples.}
\begin{senum}
\item $-5$
\item $\frac{1}{2}$
\item $\sqrt{2}$
\end{senum}

\section*{Set Builder Notation}
For each set $A$ given below, list three elements of the set. Then, write the given number $x$ in a form that shows that $x\in A$.
\emph{Note: Answers here will vary. Those given are just some possible examples.}
\begin{senum}
\item Elements: $\frac{1}{8}$, $\frac{2}{8}$, $\frac{3}{8}$. $x=\frac{4}{8}\in A$
\item Elements: $\frac{1}{2}$, $\frac{3}{8}$, $\frac{5}{6}$. $x=\frac{3}{4}\in A$
\item Elements: $1=3-2$, $2=4-2$, $3=5-2$. $x=1-2\in A$
\item Elements: $\frac{1}{2}$, $\frac{1}{7}$, $\frac{2}{3}=\frac{1}{\frac{3}{2}}$. $x=\frac{1}{\frac{1}{3}}\in A$
\end{senum}

\section*{Properties of Real Numbers/Simplifying Expressions}
Simplify each expression below. Clearly show your work on a separate sheet of paper.
\begin{smenum}
\mitemxxx
{$5x$}
{$-6a$}
{$5t-1$}
\mitemxxxx
{$2x-8$}
{$2t+12$}
{$4t+10$}
{$\frac{1}{2}x+\frac{5}{2}$}
\mitemxxxx
{$-x+30$}
{$\frac{5}{6}x+\frac{3}{2}$}
{$-3.5x + 8.5$}
{$2.3x-1.5$}
\mitemxx
{$x+y+14$}
{$25.5x-73.8$}
\end{smenum}

\section*{Properties of Equality/Solving Equations}
Solve each equation below. Clearly show your work on a separate sheet of paper.
\begin{smenum}
\mitemxx
{$x=5$}
{$x=-\frac{4}{11}$}
\mitemxx
{$y=\frac{20}{3}$}
{$y=16$}
\mitemxx
{No solution}
{$a=2$}
\mitemxx
{All real numbers are solutions}
{$a=-9$}
\mitemxx
{No solution}
{$x=-8$}
\mitemxx
{All real numbers are solutions}
{$a=-9$}
\mitemxx
{$v=4.1$}
{No solution}
\end{smenum}

\end{document}